\documentclass[11pt]{article}
\usepackage[a4paper,margin=1in]{geometry}
\usepackage{fontspec}
\usepackage[boldfont=false]{xeCJK}
\setCJKmainfont{FandolSong-Regular.otf}
\usepackage{amsmath}
\usepackage{amssymb}
\usepackage{array}
\usepackage{booktabs}
\usepackage{graphicx}
\usepackage{adjustbox}
\usepackage{enumitem}
\setlist[itemize]{topsep=2pt,partopsep=0pt,parsep=0pt,itemsep=2pt,leftmargin=1.4em}
\setlist[enumerate]{topsep=2pt,partopsep=0pt,parsep=0pt,itemsep=2pt,leftmargin=1.6em}
\usepackage{parskip}
\usepackage{fvextra}
\fvset{breaklines=true,breakanywhere=true,fontsize=\small}
\setlength{\parindent}{0pt}

\begin{document}

\begin{center}
  {\LARGE\bfseries Marketing (A Level)}\\[0.35em]
  {\large A Level Business}
\end{center}
\vspace{0.6em}

\section*{Marketing analysis}

\textbf{Marketing analysis} 营销分析 means using data to understand the market and to make better marketing decisions. At A Level you study some numbers tools: sales forecasting, time-series analysis, correlation, and elasticity.

\section*{Sales forecasting and time-series analysis}

\textbf{Sales forecasting} 销售预测 is predicting future sales from past data and market knowledge. Good forecasts help a firm plan production, staff, stock and cash.

\textbf{Time-series analysis} 时间序列分析 studies past sales recorded over time to find a pattern. It has two main parts:

\begin{itemize}
\item the \textbf{trend} 趋势 — the general direction of sales over the long run (up, down or flat).

\item the \textbf{seasonal variation} 季节性波动 — the regular rise and fall within a year (e.g. ice cream sells more in summer).

\end{itemize}

To find the trend, firms use a \textbf{moving average} 移动平均, which smooths the ups and downs by averaging several periods in a row. The seasonal variation is then how far the actual figure sits above or below the trend:

\[
\text{seasonal variation} = \text{actual value} - \text{trend value}
\]

A forecast extends the trend into the future (called \textbf{extrapolation} 外推) and then adds the usual seasonal variation. This works only if past patterns continue, so forecasts are never certain.

\section*{Correlation}

\textbf{Correlation} 相关性 measures how strongly two things move together — for example, advertising spending and sales.

\begin{itemize}
\item \textbf{positive correlation} 正相关 — when one rises, the other rises too.

\item \textbf{negative correlation} 负相关 — when one rises, the other falls.

\end{itemize}

Strong correlation helps a firm predict, but it does not prove that one thing \emph{causes} the other. Other factors may be at work.

\section*{Using elasticity in marketing decisions}

\textbf{Elasticity} 弹性 measures how much demand reacts to a change. Three types guide marketing decisions.

\begin{itemize}
\item \textbf{price elasticity of demand} 需求价格弹性 (PED) — for a \textbf{price elastic} 富有弹性 product, cutting the price raises total revenue; for a \textbf{price inelastic} 缺乏弹性 product, raising the price raises total revenue.

\item \textbf{income elasticity of demand} 需求收入弹性 (YED) — tells a firm how sales will change as incomes rise or fall, which helps plan for booms and recessions.

\item \textbf{cross elasticity of demand} 需求交叉弹性 (XED) — tells a firm how a rival's price change, or the price of a partner product, will affect its own sales.

\end{itemize}

So elasticity helps a firm set prices, choose products, and predict the effect of a competitor's move.

\section*{Marketing planning and strategy}

A \textbf{marketing strategy} 营销战略 is the long-term plan for how marketing will meet the firm's objectives. \textbf{Marketing planning} 营销规划 is the work of setting that plan: studying the market, choosing target customers, and deciding the marketing mix. Two tools shape strategy: Ansoff's matrix and product portfolio analysis.

\section*{Ansoff's matrix}

\textbf{Ansoff's matrix} 安索夫矩阵 shows four ways to grow, by mixing existing or new products with existing or new markets. The risk rises as you move away from what you know.

\begin{center}
\adjustbox{max width=\linewidth}{%
\begin{tabular}{lll}
\toprule
\textbf{} & \textbf{Existing products} & \textbf{New products} \\
\midrule
\textbf{Existing markets} & \textbf{market penetration} 市场渗透 (sell more to current buyers) & \textbf{product development} 产品开发 (new products for current buyers) \\
\textbf{New markets} & \textbf{market development} 市场开发 (current products to new buyers) & \textbf{diversification} 多元化 (new products and new markets) \\
\bottomrule
\end{tabular}}
\end{center}

Market penetration is the safest; diversification is the riskiest because both the product and the market are new.

\section*{Product portfolio analysis}

\textbf{Product portfolio analysis} 产品组合分析 looks at the firm's whole range of products together, using a tool like the Boston Matrix. It checks the balance: a firm needs some steady earners to fund its risky new products. The aim is a healthy mix, so the firm is not left with only ageing products.

\section*{Entering international markets}

To grow, many firms sell abroad. \textbf{Globalisation} 全球化 (the world becoming more connected) has made this easier. A firm can enter a new country by \textbf{exporting} 出口, by setting up a \textbf{joint venture} 合资企业 with a local partner, or by opening its own branch there.

But a strategy that works at home may need changes. \textbf{Localisation} 本地化 means adapting the marketing mix to the new market — for example, changing the product, the language, the price, or the way it is sold to suit local tastes, laws and incomes.


\end{document}
